\documentclass[11pt,a4paper]{article}
%% Shared SCRS preamble fragments (input after \documentclass)
\usepackage[margin=2.5cm]{geometry}
\usepackage{setspace}
\usepackage{authblk}
\usepackage{graphicx}
\usepackage{booktabs}
\usepackage{array}
\usepackage{placeins}
\usepackage{float}
\usepackage{amsmath,amssymb}
\usepackage{enumitem}
\usepackage{url}
\usepackage[round]{natbib}
\usepackage{hyperref}
\hypersetup{colorlinks=true,linkcolor=black,citecolor=black,urlcolor=blue}

\setstretch{1.0}
\parskip=6pt
\parindent=0pt
\pagestyle{empty}
\emergencystretch=3em
\graphicspath{{figs/}}

\newcommand{\incfig}[2][]{%
  \IfFileExists{figs/#2}{%
    \includegraphics[#1]{#2}%
  }{%
    \fbox{\begin{minipage}{0.92\linewidth}\centering\vspace{1.5em}
      Missing figure: \texttt{#2}\\[0.5em]
      {\small regenerate from \texttt{Rmd/sma-uncertainty.Rmd}}
      \vspace{1.5em}\end{minipage}}%
  }%
}

% generated by scripts/refresh-paper-numbers.R; do not edit
\newcommand{\nDraw}{1000}
\newcommand{\jabbaRun}{JABBA3}
\newcommand{\histEnd}{2024}
\newcommand{\linN}{20{,}000}
\newcommand{\linCells}{60}
\newcommand{\linMaxDiff}{4}
\newcommand{\linMeanDiff}{1.0}
\newcommand{\linWithinOne}{45}
\newcommand{\linMaxRounded}{4}
\newcommand{\trDiff}{3.3}
\newcommand{\statusGreenMvln}{52}
\newcommand{\statusGreenMcmc}{59}
\newcommand{\statusGreenPe}{79}
\newcommand{\statusGreenJabba}{76}
\newcommand{\statusGreenJabbaDyn}{86}
\newcommand{\histBMvln}{1.01 (0.80--1.30)}
\newcommand{\histFMvln}{0.73 (0.51--1.08)}
\newcommand{\histBMcmc}{1.11 (0.84--1.46)}
\newcommand{\histFMcmc}{0.96 (0.67--1.35)}
\newcommand{\histBMvlnMc}{n/a}
\newcommand{\histFMvlnMc}{n/a}
\newcommand{\histBJabba}{1.22 (0.67--2.18)}
\newcommand{\histFJabba}{0.29 (0.07--1.19)}
\newcommand{\histBJabbaDyn}{1.31 (0.82--2.14)}
\newcommand{\histFJabbaDyn}{0.22 (0.05--0.87)}
\newcommand{\histBJabbaOneB}{1.31 (0.82--2.14)}
\newcommand{\histFJabbaOneB}{0.22 (0.05--0.87)}
\newcommand{\histBJabbaOneS}{1.22 (0.67--2.18)}
\newcommand{\histFJabbaOneS}{0.29 (0.07--1.19)}
\newcommand{\histBJabbaTwoS}{n/a}
\newcommand{\histFJabbaTwoS}{n/a}
\newcommand{\greenZeroMvlnFortyFive}{96}
\newcommand{\tacMvlnFortyFiveSixty}{1,288}
\newcommand{\tacMvlnFortyFiveSeventy}{1,033}
\newcommand{\greenZeroMvlnSeventy}{100}
\newcommand{\tacMvlnSeventySixty}{1,587}
\newcommand{\tacMvlnSeventySeventy}{1,486}
\newcommand{\greenZeroMcmcFortyFive}{99}
\newcommand{\tacMcmcFortyFiveSixty}{1,185}
\newcommand{\tacMcmcFortyFiveSeventy}{1,080}
\newcommand{\greenZeroMcmcSeventy}{100}
\newcommand{\tacMcmcSeventySixty}{1,188}
\newcommand{\tacMcmcSeventySeventy}{1,114}
\newcommand{\greenZeroPeFortyFive}{100}
\newcommand{\tacPeFortyFiveSixty}{1,320}
\newcommand{\tacPeFortyFiveSeventy}{1,150}
\newcommand{\greenZeroPeSeventy}{100}
\newcommand{\tacPeSeventySixty}{1,300}
\newcommand{\tacPeSeventySeventy}{1,217}
\newcommand{\greenZeroJabbaFortyFive}{97}
\newcommand{\tacJabbaFortyFiveSixty}{2,300}
\newcommand{\tacJabbaFortyFiveSeventy}{2,300}
\newcommand{\greenZeroJabbaSeventy}{98}
\newcommand{\tacJabbaSeventySixty}{2,300}
\newcommand{\tacJabbaSeventySeventy}{2,300}
\newcommand{\greenZeroJabbaDynFortyFive}{98}
\newcommand{\tacJabbaDynFortyFiveSixty}{2,300}
\newcommand{\tacJabbaDynFortyFiveSeventy}{2,300}
\newcommand{\greenZeroJabbaDynSeventy}{100}
\newcommand{\tacJabbaDynSeventySixty}{2,300}
\newcommand{\tacJabbaDynSeventySeventy}{2,300}
\newcommand{\nMcmc}{1,485}
\newcommand{\mcmcCatches}{26}
\newcommand{\apxMeanMvln}{n/a}
\newcommand{\apxMaxMvln}{n/a}
\newcommand{\apxMeanMcmc}{n/a}
\newcommand{\apxMaxMcmc}{n/a}
\newcommand{\apxMeanJabba}{n/a}
\newcommand{\apxMaxJabba}{n/a}
\newcommand{\tacRangeSixtySeventy}{713}
\newcommand{\rpBaseSsbmsy}{552}
\newcommand{\rpBaseMsy}{1,880}
\newcommand{\rpMleSsbmsy}{360}
\newcommand{\rpMleMsy}{1,240}
\newcommand{\rpMcSsbmsy}{369}
\newcommand{\rpMcMsy}{1,280}
\newcommand{\rpSsbmsyDiff}{33}
\newcommand{\rpMsyDiff}{32}
\newcommand{\rpFDiff}{29}

\InputIfFileExists{numbers-ssproject.tex}{}{}

\newcommand{\PG}{P(\mathrm{green})}
\newcommand{\BB}{B/B_{\mathrm{MSY}}}
\newcommand{\SB}{\mathrm{SSF}/\mathrm{SSF_{MSY}}}
\newcommand{\FF}{F/F_{\mathrm{MSY}}}

\date{}

\renewcommand\thesection{\arabic{section}.}
\renewcommand\thesubsection{\arabic{section}.\arabic{subsection}}
\makeatletter
\renewcommand\section{\@startsection{section}{1}{\z@}%
  {-1.5ex \@plus -1ex \@minus -.2ex}{0.8ex \@plus .2ex}{\normalfont\bfseries}}
\renewcommand\subsection{\@startsection{subsection}{2}{\z@}%
  {-1.25ex \@plus -1ex \@minus -.2ex}{0.6ex \@plus .2ex}{\normalfont\itshape}}
\makeatother


\title{\textbf{Outlines for papers not yet written}\\[0.5em]
\large Not an SCRS submission. The projection paper and the journal draft
already exist.}

\author[1]{L. T. Kell}
\affil[1]{Centre for Environmental Policy, Imperial College London, London, United Kingdom.\\
\texttt{laurie@seaplusplus.co.uk}}

\begin{document}
\maketitle

\begin{center}\textit{WHAT ALREADY EXISTS}\end{center}

Two manuscripts are in progress. Neither is outlined again here.

\begin{itemize}[leftmargin=*]
  \item \textbf{Journal draft,} \texttt{peer\_review\_draft.tex}.
    Rebuilding-time reference points for non-retained bycatch, with North
    Atlantic shortfin mako as the example. Argument, status-versus-projection
    framing, and the case that the target need not be $B_{\mathrm{MSY}}$.
  \item \textbf{SCRS paper,} \texttt{scrs\_2026\_proj\_supplement.tex}.
    Submitted to ICCAT with its introduction and discussion around the four
    projections. After withdrawal, that file is the supplementary material
    of the journal manuscript. It is not rewritten for the journal.
\end{itemize}

The WGSAM prospectus and the three-way split (papers~i--iii) are in
\texttt{archive/}. The closed-form check against the published Kobe matrix is
Appendix~A of the projection paper. Two pieces in the archived drafts are
not in either current manuscript. The weighting paper is written after the
submitted ensemble manuscript. The process-error note can be written from
the cached projection. Neither is ready to number as an SCRS document.

\section{Selecting, excluding and weighting projection runs}

\textbf{Possible form.} A later SCRS paper for WGSAM, written after
\citet{WangEnsembleR5}. That manuscript is in revision as ICESJMS-2026-566
(Manuscript-R6, \path{C:/work/papers/02_submitted/ensemble}). It is the
weighting method: diagnostics first, then predictive weights on an already
credible set. This paper applies that method to projection runs. The Shark
Species Group asked for an objective way to select, exclude and weight
models across frameworks \citep{SharkWG2026}. Do not number this paper until
that manuscript is the method being cited, and until the SMA-N weights have
been computed with it.

\textbf{Claim.} A Kobe matrix that averages SS3 and JABBA3 without a gate
treats a structural disagreement as a weight. The projection paper shows the
disagreement: at $\PG = 0.6$ in 2070 the SS3-MVLN catch is
\tacMvlnSeventySixty~t and both JABBA3 catches lie above 2{,}300~t. Weighting
is the decision about whether that second tier enters the advice table.

\textbf{Audience.} WGSAM. SMA-N is the sketch, not a recommended catch.

\textbf{Sections.}
\begin{enumerate}[leftmargin=*]
  \item Why an unweighted mix of frameworks is not an ensemble. The method
    is the one in \citet{WangEnsembleR5}, not a scheme invented for SMA-N.
  \item Five steps, in this order.
    \begin{itemize}[leftmargin=*]
      \item \textbf{Select.} Keep runs that pass agreed diagnostics: fits,
        retrospectives or hindcasts, residual structure
        \citep{Kell2026Diags,Carvalho2021}.
      \item \textbf{Exclude.} Drop runs that fail those diagnostics, and drop
        a projection whose biology cannot answer the question. For SMA-N that
        is a surplus-production projection when recovery depends on the lag
        between exploitable biomass and spawners \citep{Winker2020SMA}.
      \item \textbf{Propagate.} Closed-form MVLN when the delta-method model
        is log-normal. MCMC when the posterior is not, or when parameters are
        bounded. Recruitment process error when the advice is a multi-year
        recovery and year-by-year MVLN would understate risk near the cliff.
      \item \textbf{Weight.} Only among runs that remain, using the predictive
        weights of \citet{WangEnsembleR5}. That paper weights an already
        credible set. It does not decide whether a surplus-production
        projection is eligible for an SMA-N rebuilding matrix. Excluded runs
        stay out of the average.
      \item \textbf{Report.} Show MVLN, MCMC and any process-error sensitivity
        as their own columns, even if a weighted mean is also shown. Keep that
        table separate from Species Group catch advice.
    \end{itemize}
  \item Candidate runs for the sketch: SS3-MVLN, SS3-MCMC (only after the
    same \texttt{do\_recdev} as MVLN), FLasher process error as a sensitivity,
    the JABBA3 bridge, and the WG JAGS runs if they are ported. The JAGS 1-B
    and 1-S results from an older knit are not relabelled as JABBA3.
  \item Worked weights: diagnostic outcome for each run, who is excluded, the
    weight scheme, and $\PG$ and the catch-at-target with and without the
    weights. Reuse the projection paper's catch table as the unweighted
    baseline.
  \item Decision the paper has to make: does the JABBA3 tier enter the SMA-N
    advice table, and if so with what weight? The projection paper shows that
    the tier changes the catch. It does not make that decision.
\end{enumerate}

\textbf{Already written.} The weighting method, in the submitted ensemble
manuscript. The five-step box for this application is in
\texttt{archive/scrs\_2026\_ensemble\_weights.tex}. The candidate list is a
sketch. The gate column is empty.

\textbf{Missing before it is a paper.} Cite \citet{WangEnsembleR5} in the
form that will be the method, then a diagnostic pass or fail for each SMA-N
run, a weight from that scheme for each retained run, and one table of
weighted against unweighted $\PG$ and catch-at-target.

\section{Recruitment process error at the Kobe cliff}

\textbf{Possible form.} Short SCRS sensitivity note. It is not one of the
four projections in the supplement. SS3-MVLN has no recruitment path: each
year's $(\log B, \log F)$ is independent. For a late-maturing shark the
advice is a multi-year recovery, so that omission is a methods question of
its own.

\textbf{Claim.} Given the same MCMC terminal state, carrying historical
recruitment residuals forward changes the catch that meets $\PG = 0.6$ near
the steep part of the risk curve. Spawning biomass stays smooth, because many
year classes filter a single residual. The risk does not.

\textbf{Audience.} WGSAM. The numerical example is SMA-N, conditioned on the
recDev2 MCMC end state, projected with \texttt{FLRebuild::ssProject} and
\texttt{FLasher::fwd}. Residuals are circular-block resamples of the
historical series, not independent draws by year
(\texttt{scripts/run-ss-project.R}, \texttt{data/runs/ss-project.rds}).

\textbf{Sections.}
\begin{enumerate}[leftmargin=*]
  \item What MVLN leaves out: no recruitment trajectory, so a late-maturing
    stock's recovery is a sequence of independent margins.
  \item The contrast that isolates process error. Both paths start from the
    MCMC terminal state. One uses residual $= 1$. The other resamples the
    historical residual series in blocks. This is not the MVLN-versus-MCMC
    comparison, which is still confounded by \texttt{do\_recdev}.
  \item Trajectories: median and 90\% interval of $\BB$ and $\FF$ at a few
    catches (0, 1{,}200, 1{,}883, 2{,}300~t). Figure
    \texttt{figs/ss-project-pe.png}.
  \item $\PG$ by catch and year for the deterministic path and the resampled
    path (\texttt{tabs/ss-project.tex}). The catch at 60\% and 70\% in 2045
    and 2070: \tacPeFortyFiveSixty~t and \tacPeSeventySixty~t at 60\%, and
    \tacPeFortyFiveSeventy~t and \tacPeSeventySeventy~t at 70\%.
  \item Where it matters. Away from the cliff the two paths agree. On the
    steep part of the curve, process error turns a step in $\PG$ into a
    slope and fattens the upper tail of $\FF$. Age structure filters the
    biomass path and does not filter that risk
    \citep{Winker2020SMA}.
  \item What the note will not claim. Absolute $\PG$ is an FLR projection of
    the SS3 end state. The result is the difference between residual $= 1$
    and the block resample, not a replacement for the endorsed SS3 matrix
    \citep{Liniers2026}.
\end{enumerate}

\textbf{Already computed.} The cached \texttt{ssProject} run, the figure, and
the catch numbers above.

\textbf{Missing before it is a paper.} A same-\texttt{do\_recdev} MVLN next
to this sensitivity, so the process-error effect is not only shown against
the deterministic FLasher path. A sentence on how large a change in catch
would alter Rec.~21-09 advice.

\bibliographystyle{plainnat}
\bibliography{refs}

\end{document}
